% Research note, 2026-09-30. Independent of main.tex.
\documentclass[11pt]{article}
\usepackage{amsmath,amssymb,amsthm,hyperref}
\newtheorem{theorem}{Theorem}
\newtheorem{lemma}[theorem]{Lemma}
\newtheorem{corollary}[theorem]{Corollary}
\newcommand{\vii}{v_2}
\newcommand{\Zii}{\mathbb Z_2}
\newcommand{\Qii}{\mathbb Q_2}
\begin{document}
\title{An exponential obstruction for rational equal-weight interval quadrature}
\author{Research note; proof candidate for independent verification}
\date{30 September 2026}
\maketitle

\paragraph{Historical attribution.}
The parity polynomial $\sum_{j=1}^r(-2)^{j-1}\binom{x}{j}$ used below is
classical: Mit'kin \cite[Lemma~7 and following remark]{mitkin} credits it
to Hua (1940), and records related Hilbert--Kamke applications by Arkhipov.
The present note studies its integral after dilation and the consequence
for rational equal-weight interval quadrature; no novelty claim is made
for the polynomial itself.


\begin{theorem}\label{thm:main}
Let $x_1,\ldots,x_K$ be rational numbers (not necessarily in $[0,1]$)
with
\[
 \frac1K\sum_{i=1}^K x_i^j=\frac1{j+1}\qquad(0\leq j\leq r).
\]
Then $K\geq 2^{\lfloor r/2\rfloor}$. More precisely:
\begin{enumerate}
\item If every $x_i$ belongs to $\Zii$, then $2^r\mid K$.
\item Otherwise, after the normalization below, the nonzero number of odd nodes is divisible by $2^{\min(r,\vii(K)+\lfloor r/2\rfloor)}$.
\end{enumerate}
Repeated nodes are allowed.
\end{theorem}

\paragraph{Status.}
The following is a complete proposed proof, with its coefficient identities
also checked by exact rational arithmetic in the accompanying script.
It is not a proof of the per-die exponential conjecture: the insertion
reduction in the manuscript is only one-way. If valid, it rules out the
proposed subexponential rational-quadrature route to counterexamples.

Define $G(z)=z/\log(1+z)$ as a formal series over $\mathbb Q$, with constant
term one. The coefficient $[z^j]G(z)$ is
$\int_0^1\binom{x}{j}\,dx$, as follows by integrating the formal identity
$(1+z)^x=\sum_{j\geq0}\binom{x}{j}z^j$.

Put
\[
 D(t)=\frac{-\log(1-2t)}{2t}
     =\sum_{j\geq0}\frac{2^j}{j+1}t^j.
\]
Its coefficients are in $\Zii$, and $D(t)\equiv1+t\pmod2$: the terms
$j=0,1$ are units, whereas $j-\vii(j+1)\geq1$ for $j\geq2$.
Consequently $G(-2t)=D(t)^{-1}\equiv(1+t)^{-1}\pmod2$, so every
coefficient of $G(-2t)$ is odd. In particular,
\begin{equation}\label{eq:Gregory}
 \vii\left(\int_0^1\binom{x}{j}\,dx\right)=-j.
\end{equation}
If every $x_i\in\Zii$, then $\binom{x_i}{r}\in\Zii$.
Exactness and \eqref{eq:Gregory} give $\vii(K)\geq r$, proving the first case.

For the other case, consider the degree-$r$ polynomial
\[
 H_r(y)=\sum_{j=1}^r(-2)^{j-1}\binom{y}{j}.
\]
For every $y\in\Zii$,
\begin{equation}\label{eq:parity}
 H_r(y)\equiv\boldsymbol1_{y\notin2\Zii}\pmod{2^r}.
\end{equation}
Indeed, for a nonnegative integer $y$, the binomial expansion of $(-1)^y$
gives this congruence after discarding terms of index greater than $r$.
Continuity and density of the nonnegative integers extend it to $\Zii$.

\begin{lemma}\label{lem:integral}
For every integer $a\geq1$ and every integer $r\geq1$,
\[
 \vii\left(\int_0^1H_r(2^a x)\,dx\right)\geq\lfloor r/2\rfloor.
\]
\end{lemma}
\begin{proof}
Write $G(z)=\sum_{j\geq0}g_jz^j$. We proved above that
$\vii(g_j)=-j$. For an integer $N\geq1$, coefficientwise integration gives
\[
 F_N(z):=\int_0^1(1+z)^{Nx}\,dx
   =\frac{(1+z)^N-1}{N\log(1+z)}
   =G\bigl((1+z)^N-1\bigr).
\]
The last identity follows from the formal identity
$\log((1+z)^N)=N\log(1+z)$.

Suppose $N$ is even and set $W_N(t)=(1-2t)^N-1$.
Its coefficient of degree $k\geq1$ is $(-2)^k\binom Nk$.
For $k=1$ this has valuation at least $2$; for $k\geq2$ its valuation
is at least $k\geq\lceil k/2\rceil+1$. Hence every coefficient of
$W_N(t)^j$ in degree $k$ has valuation at least $\lceil k/2\rceil+j$.
Multiplying by $g_j$ and summing proves that
\[
 F_N(-2t)=G(W_N(t))=\sum_{k\geq0} A_k t^k,
 \qquad A_0=1,\qquad \vii(A_k)\geq\lceil k/2\rceil\quad(k\geq1).
\]
All these series converge on $\Zii$. More precisely,
$W_N(t)\in4\Zii$ for $t\in\Zii$, so the valuation of
$g_jW_N(t)^j$ is at least $j$, uniformly in $t$. Evaluation at $t=1$
therefore commutes with the composition, and
\[
 \sum_{k\geq0}A_k=G(W_N(1))=G(0)=1.
\]
Apply this with $N=2^a$.
Finally,
\[
 \int_0^1 H_r(2^a x)\,dx
  =-\frac12\sum_{k=1}^r A_k
  =\frac12\sum_{k=r+1}^\infty A_k.
\]
Its valuation is at least $\lceil(r+1)/2\rceil-1=\lfloor r/2\rfloor$.
\end{proof}

\begin{proof}[Conclusion of Theorem~\ref{thm:main}]
Suppose some $x_i$ is not $2$-adically integral. Choose the smallest
positive integer $a$ such that $y_i=2^a x_i\in\Zii$ for every $i$.
At least one $y_i$ is odd. Let $q$ be the number of odd $y_i$, so
$1\leq q\leq K$. Exactness, \eqref{eq:parity}, and
Lemma~\ref{lem:integral} yield
\[
 q\equiv\sum_{i=1}^K H_r(y_i)
   =K\int_0^1 H_r(2^a x)\,dx\pmod{2^r}.
\]
In particular, $2^{\lfloor r/2\rfloor}\mid q$, proving $K\geq2^{\lfloor r/2\rfloor}$.
The displayed congruence also gives the more precise divisibility in item (2).
\end{proof}

\begin{thebibliography}{9}
\bibitem{mitkin} D.~A.~Mit'kin, \emph{An estimate for the number of terms in the Hilbert--Kamke problem}, Math. USSR Sbornik \textbf{57} (1987), no.~2, 561--590, Lemma~7 and its following remark, \href{https://www.mathnet.ru/eng/sm1845}{primary journal record and full text}.
\end{thebibliography}
\end{document}
